% ============================================================
%  Temario — Estadística I
%  Licenciatura en Economía, CIDE. Otoño 2026.
% ============================================================

\documentclass[letterpaper, 12pt]{article}
\pagestyle{plain}

\usepackage[utf8]{inputenc}     % Acentos
\usepackage[T1]{fontenc}
\usepackage[spanish]{babel}     % Idioma español

\usepackage{amsfonts}           % Matemáticas
\usepackage{amsmath}
\usepackage{amssymb}

\usepackage{graphicx}           % Gráficos e imágenes
\usepackage[margin=2.5cm]{geometry}

\usepackage[colorlinks=true, linkcolor=black, urlcolor=blue,
            citecolor=blue]{hyperref}

\title{\textbf{Centro de Investigación y Docencia Económicas, A.C.} \\ División de Economía \\ \textbf{Estadística I}. \\ Programa de Licenciatura en Economía \\ Otoño 2026}

\author{Benjamín Oliva-Vázquez \\ (\url{benjamin.oliva@cide.edu} \\ y \\ 
 \url{benjov@ciencias.unam.mx}) \\ \\ \textbf{Laboratorista}: \\ Giovanni Paulo Ordonez Corona \\ (\url{giovanni.ordonez@alumnos.cide.edu}) }

\date{Agosto 2026}

\begin{document}
\maketitle

\textbf{Objetivos}: 
\begin{enumerate}
    \item Desarrollar las herramientas teóricas y prácticas de probabilidad e introducir el análisis de estadística descriptiva e inferencia estadística.
    
    \item Establecer las bases para el uso de paquetes estadísticos y de procesamiento de datos en R.
\end{enumerate}

\textbf{Materiales}: el temario, los laboratorios y el código de clase están disponibles en \href{https://github.com/benjov/Estadistica-I-2026}{GitHub}.\footnote{\url{https://github.com/benjov/Estadistica-I-2026}} La bibliografía está disponible, sólo para las personas integrantes del grupo, en \href{https://drive.google.com/drive/folders/1eQlLrz9h-DBz7NLsZazcqHM0bCM204w_?usp=sharing}{Google Drive}.

\textbf{Horarios:} 
\begin{enumerate}
\item Clases: viernes de 9:00 a 12:00 hrs en el Aula de Videoconferencias del edificio Santa Fe.

\item Las horas de oficina podrán ser en línea o presenciales y se programarán en función de los requerimientos de cada persona, y 

\item Laboratorios: lunes a las 9:40 hrs en el mismo salón (Aula de Videoconferencias, edificio Santa Fe).
\end{enumerate}

\section{Temario}

\begin{enumerate}
\item Probabilidad
	\begin{enumerate}
    	\item Espacios Muestrales y Eventos
        \item Cálculo Combinatorio
        \item Definición de Probabilidad
        \item Independencia
        \item Probabilidad Condicional
        \item Teorema de Bayes
    \end{enumerate}
\item Variables Aleatorias
	\begin{enumerate}
    	\item Funciones de Distribución y de Probabilidad
        \item Algunas funciones de Variables Aleatorias Discretas
        \item Algunas funciones de Variables Aleatorias Continuas
        \item Distribuciones Bivariadas
        \item Distribuciones Marginales
        \item Variables Aleatorias Independientes
        \item Distribuciones Condicionales 
        \item Distribuciones Multivariadas y Muestras IID
        \item Transformaciones\footnote{Tema optativo y se verá según disponibilidad de tiempo.}
    \end{enumerate}
\item Esperanza
	\begin{enumerate}
    	\item Esperanza de variables aleatorias
        \item Propiedades de la Esperanza
        \item Varianza y Covarianza
        \item Esperanza y Varianza de algunas variables aleatorias
        \item Esperanza Condicional
        \item Funciones Generadoras de Momentos
    \end{enumerate}
\item Algunas Desigualdades Importantes
    \begin{enumerate}
        \item Desigualdades de Probabilidad
        \item Desigualdades de Esperanza
    \end{enumerate}
\item Convergencia de Variables Aleatorias
    \begin{enumerate}
        \item Tipos de convergencia
        \item La Ley de los Grandes Números
        \item El Teorema de Límite Central
    \end{enumerate}
\item Estadística Descriptiva con R (tema transversal que se verá durante la primera mitad del curso)
	\begin{enumerate}
    	\item Conceptos elementales
        \item Descripciones numéricas
        \item Descripciones gráficas
        \item Descripciones para datos conjuntos
    \end{enumerate}
\end{enumerate}

\section{Bibliografía}
\begin{itemize}
    \item Agresti, Alan y Maria Kateri (2021) \textit{Foundations of Statistics for Data Scientists: With R and Python}. Ed. Chapman and Hall/CRC.
    
    \item Larsen, Richard J., y Morris L. Marx (2018) \textit{An Introduction to Mathematical Statistics and its Applications}. 6ta Edición. Ed. Pearson. (*) 
    
    \item Miller, Irwin, y Marylees Miller (2014) \textit{John E. Freund's Mathematical Statistics with Applications}. 8va. Edición. Ed. Pearson. (*)
    
    \item Wackerly, Dennis D., William Mendenhall III y Richard L. Scheaffer (2008) \textit{Mathematical Statistics with Applications}. 7ma Edición. Ed. Thomson Learning. (*)
    
    \item Wasserman, Larry (2004) \textit{All of Statistics. A Concise Course in Statistical Inference}. Springer Texts in Statistics. Springer. 
\end{itemize}

Todos los textos estarán disponibles sólo para las personas integrantes del grupo en \href{https://drive.google.com/drive/folders/1eQlLrz9h-DBz7NLsZazcqHM0bCM204w_?usp=sharing}{Google Drive}.

Notas: (*) Textos base o principales del curso.

\section{Evaluación}
\begin{enumerate}
    \item Tareas, quizzes y otras asignaciones en R (15\%)
    \item Laboratorio (15\%)
    \item Examen Parcial (35\%)
    \item Examen Final (35\%)
\end{enumerate}

\section{Calendario (tentativo)}

\textbf{Nota}: Las fechas mostradas a continuación están sujetas a cambios en función de un acuerdo o consenso del grupo.

\begin{table}[htb]
\centering
\begin{tabular}{| c | c |}
\hline
\hline
    \textbf{Laboratorio / Examen} & \textbf{Fecha de Entrega o} \\
     & \textbf{Aplicación (tentativa)} \\
\hline \hline
    Laboratorio 1 & 31 de agosto de 2026 \\ \hline
    Laboratorio 2 & 7 de septiembre de 2026 \\ \hline
    Laboratorio 3 & 14 de septiembre de 2026 \\ \hline
    Laboratorio 4 & 21 de septiembre de 2026 \\ \hline
\hline
    Examen Parcial (temas 1 y 2*) & 2 de octubre de 2026 \\ \hline
\hline
    Laboratorio 5 & 19 de octubre de 2026 \\ \hline
    Laboratorio 6 & 26 de octubre de 2026 \\ \hline
    Laboratorio 7 & 2 de noviembre de 2026 \\ \hline
    Laboratorio 8 & 9 de noviembre de 2026 \\ \hline
\hline
    Examen Final (temas 2*, 3, 4 y 5) & 11 de diciembre de 2026 \\ \hline
\hline
\multicolumn{2}{l}{\textit{Notas}: * Este tema podría evaluarse de forma parcial} \\
\multicolumn{2}{l}{dependiendo del avance del curso.}\\
\end{tabular}
\label{Calendario}
\end{table}

\end{document}
